\section{Questions left open}
The results leave several different algorithmic problems. Each depends on
the written presentation, so the following questions keep the input and
the desired output explicit.

\paragraph{Infinite height and a varying height parameter.}
For fixed $p$, every fixed exact finite height and every fixed lower
threshold at least two is $\ModP$-complete on the smooth anticanonical
circuit promise. Every finite height has an explicit prime-field seed.
Theorem~\ref{fg:thm:space} computes the full height, including infinity,
in polynomial space. What is the complexity of deciding infinite
height? The fixed-threshold reductions preserve only finitely many
logarithm coefficients and do not give a lower bound for this predicate.
A lower-bound construction must enforce a genuine infinite-height
alternative or realize an appropriate computation in the geometric
finite operator. When the finite queried height or threshold is given
in binary, the decision problems are in polynomial space and are
$\ModP$-hard by restriction to a fixed query. A matching classification
in that regime is also not proved. In fixed dimension the full height
is polynomial-time computable, whereas for a growing-dimensional sparse
equation the proved polynomial bounds concern fixed finite queries.

\paragraph{Unrestricted saturation.}
Fix a prime and a codimension. Given a homogeneous sparse list with binary
exponents whose saturation is a smooth geometrically integral
anticanonical complete-intersection ideal, can its coherent Frobenius
scalar be computed in polynomial time? For a hypersurface, the remaining
case lies beyond the small-cofactor and polynomial degree-and-support
conditions, even after the safe preprocessing of
Lemma~\ref{sp:lem:preprocessing} and the three prime-field relation
theorems. The zero-sum strict family in
Example~\ref{sp:ex:zero-sum-family} moves into the polynomial
degree-and-support cofactor class under these reductions.
In higher codimension,
Theorem~\ref{sci:thm:random-excess} settles bounded complete-intersection
excess in the input ideal after scalar extension.
Theorem~\ref{sci:thm:certified-preprocessing} further permits coordinate
divisions certified from the original equations, radical-safe digit
reduction, and the prime-field relations of
Theorems~\ref{sp:thm:linearized-preprocessing},
\ref{sp:thm:affine-preprocessing}, and~\ref{sp:thm:zero-sum-preprocessing}.
This includes lists whose
least excess remains unbounded after coordinate and digit reduction
alone. The unresolved case allows the least excess of the
processed ideal to grow without a bound. The degree of the
unknown annihilator then grows, and the present linear systems can have
exponentially many coordinates. Small generators of the saturated ideal
do not bound this parameter for the original list.

\paragraph{Deterministic recovery at bounded excess.}
Under the existential bounded-excess promise alone, can a certified good
projection centre be chosen deterministically for one of the polynomially
many candidate tuples? Lemma~\ref{sci:lem:profile-candidates} constructs
the candidates deterministically, and Proposition~\ref{sci:prop:fibre-grid}
recovers the annihilator from a fixed interpolation grid once a centre
is known. The finite-field constructions of
Theorem~\ref{sci:thm:balanced} give deterministic scalar recovery under a maximum-degree
bound, and Theorem~\ref{sci:thm:gap} gives deterministic scalar recovery under a fixed degree gap.
Theorem~\ref{sci:thm:linear-prefix} also constructs a centre
and recovers the scalar deterministically when
\(\dim (K^\circ)_1=c-1\), where $K^\circ$ is the ideal generated
by the processed list. The factored-prefix theorem constructs a centre
when an indexed regular tuple has its first $c-1$ forms given as
products of linear factors, even if $(K^\circ)_1=0$.
After certified preprocessing, choosing a centre is still the only
random step. The remaining case has no such designated factored
prefix and has \(\dim (K^\circ)_1\leq c-2\).
The general bounded-excess theorem has an always-correct algorithm of
expected polynomial time. Its successful outputs are certified; the
remaining issue is deterministic selection, rather than correctness of
the returned scalar.

\paragraph{The intermediate support range.}
For a fixed odd prime and a fixed number of nonlinear equations, what is
the complexity when the anticanonical support excess
\(K=(p-1)\sum_i\sigma_i-m\) grows faster than \(\log L\) but more
slowly than every fixed positive power of the degree? The weighted
syndrome algorithm is polynomial for \(K=O(\log L)\), whereas the
smooth hard family can have \(K\leq d^\eta\) for every fixed
\(\eta>0\). These bounds do not classify, for example, a general
polylogarithmic excess. The question concerns the explicit sparse
monomials in the displayed ambient coordinates after discarding every
monomial with an individual coordinate exponent at least \(p\).

\paragraph{Effective compression over the input field.}
Can the ordinary-projective trinomial presentation be replaced by a
polynomial-size list defining the same scheme in a smaller projective
space over the original finite field? A satisfactory construction must
produce the rational centre or another embedding, its defining equations,
their degrees and monomial counts, and a bound on the construction time.
The extension and elimination bounds in Section~\ref{sp:sec:compression}
do not give that conclusion. Picard and secant arguments restrict the
possible embeddings but do not construct a short equation list. A
uniform compression of the hard family to a fixed number of explicitly
sparse equations would imply \(\ModP\subseteq\mathrm P\), because
its point-count residue could then be evaluated by the sparse algorithm.
Existence of short coordinates without an efficient way to find them is
a different question.

\paragraph{A different bounded-codimension realization.}
Could a different smooth bounded-codimension projective scheme realize
the same hard scalar with a polynomial-size explicit equation list?
Such a construction must identify its coherent group and its readout,
as well as the number of equations. Fixed codimension does not imply a
fixed number of written equations. For anticanonical global complete
intersections, generation before saturation and the bounded-excess
classes already have the algorithms proved above. Constructions outside
those classes, or with a different specified coherent group, remain
unsettled.

\paragraph{Recovering a hidden counting presentation.}
The selector and register algorithms require a marked affine Boolean
domain, its wiring or Walsh data, and a verified relation between legal
words and primary assignments. The split-quadric theorem recovers private
coordinates from a designated complete quadratic span and a marked base.
It does not extract that span and the base map from an arbitrary ideal
presentation, nor does it control the support of every further equation
after recovery. Those extraction and support costs remain part of the
problem. General private permutations need not preserve the bounded row
rank. A fixed-rank cohomological realization of the empty-fibre hardness
would also require a different mechanism: a fixed population cap in the
displayed sparse Walsh model already permits the polynomial emptiness
test.

\paragraph{The least total Betti number.}
Over a fixed finite field, how small can total geometric Betti number be
in a polynomial-size smooth projective encoding of exact witness counts,
or of their residues modulo a fixed prime? Exact decoding from the zeta
function alone cannot have uniformly bounded total Betti number. The
quantitative lower bounds are
\(B_n> (n/(6\log_2q))^{1/3}\), and
\(B_n=\Omega(\sqrt n)\) at fixed variety dimension and fixed field.
The least possible growth remains undetermined. For modular counts,
even a constant Betti bound remains open: a finite set of possible zeta
functions does not provide an efficient classifier for succinct input
equations. Arithmetic circuits, sparse lists with binary exponents, and
succinct curve function fields impose different requirements on an
encoding and must be distinguished.

\paragraph{Which efficiently computable invariant could suffice?}
The invariant-collision examples settle the proposed coarse invariants
on the stated families: Hilbert polynomials, Chern numbers, Betti numbers,
coherent dimensions, and in the varying-field curve example the entire
based coherent Frobenius operator lose the relevant residue information.
They do not classify every possible functorial invariant with a
polynomial-time algorithm. Any further proposal must specify both its
algorithm on the chosen presentation and a proved decoding map, including
the cost of decoding. A polynomial-time invariant together with a
polynomial-time decoder for the hard family's residue would imply
\(\ModP\subseteq\mathrm P\). An unrestricted decoding map makes a
different claim. Weight
bounds alone do not supply the missing multiplicities, and a virtual
class alone does not identify an actual cohomological representation.
